\documentclass{article}
\usepackage{amsmath,amssymb,amsthm}
\usepackage{algorithm}
\usepackage{algorithmic}
\usepackage{bm}
\usepackage{hyperref}
\newtheorem{theorem}{Theorem}
\newtheorem{lemma}{Lemma}
\newtheorem{assumption}{Assumption}
\newcommand{\EE}{\mathbb{E}}
\newcommand{\RR}{\mathbb{R}}
\newcommand{\norm}[1]{\left\|#1\right\|}
\newcommand{\ip}[2]{\left\langle #1,#2\right\rangle}
\begin{document}

\section*{Algorithm Name}
\textbf{Stochastic (Subsampled) Newton Method for Non‑Convex Smooth Optimization}

\section*{Mathematical Formulation}
We consider the finite‑sum (or expectation) objective
\[
\min_{x\in\RR^{d}} f(x)\;,\qquad 
f(x)=\frac{1}{n}\sum_{i=1}^{n} f_i(x)
        \;=\;\EE_{\,\xi\sim\mathcal D}[\,f_{\xi}(x)\,] .
\]

At iteration $t$ a minibatch $\mathcal S_t\subset\{1,\dots,n\}$ of size $|\mathcal S_t|=b$ is sampled
independently of the past.  The stochastic gradient and stochastic Hessian are
\[
g_t \;:=\; \frac{1}{b}\sum_{i\in\mathcal S_t}\nabla f_i(x_t),\qquad
H_t \;:=\; \frac{1}{b}\sum_{i\in\mathcal S_t}\nabla^{2} f_i(x_t).
\]
A (small) ridge parameter $\lambda>0$ is added to guarantee invertibility:
\[
\tilde H_t \;:=\; H_t + \lambda I .
\]

With a stepsize (learning rate) $\alpha>0$ the update reads
\begin{equation}
\boxed{%
x_{t+1}=x_t-\alpha\,\tilde H_t^{-1} g_t.}
\tag{1}
\end{equation}
The algorithm stops after $T$ iterations or when a prescribed tolerance on
$\norm{\nabla f(x_t)}$ is reached.

\section*{Pseudocode}
\begin{algorithm}[H]
\caption{Stochastic Subsampled Newton}
\begin{algorithmic}[1]
\STATE \textbf{Input:} initial point $x_0\in\RR^{d}$, stepsize $\alpha>0$, ridge $\lambda>0$, minibatch size $b$, max. iterations $T$.
\FOR{$t=0,1,\dots,T-1$}
    \STATE Sample a minibatch $\mathcal S_t\subset\{1,\dots,n\}$ with $|\mathcal S_t|=b$.
    \STATE Compute stochastic gradient 
          $g_t\gets \frac1b\sum_{i\in\mathcal S_t}\nabla f_i(x_t)$.
    \STATE Compute stochastic Hessian 
          $H_t\gets \frac1b\sum_{i\in\mathcal S_t}\nabla^{2}f_i(x_t)$.
    \STATE Form regularised matrix $\tilde H_t\gets H_t+\lambda I$.
    \STATE Update $x_{t+1}\gets x_t-\alpha\,\tilde H_t^{-1} g_t$.
\ENDFOR
\STATE \textbf{Output:} $x_T$ (or the best iterate seen).
\end{algorithmic}
\end{algorithm}

\section*{Dry Run Example}
Consider the scalar quadratic $f(w)=(w-3)^2$.  Then
\[
\nabla f(w)=2(w-3),\qquad \nabla^2 f(w)=2 .
\]
We take a deterministic ``subsample’’ (so $H_t=2$, $g_t=2(w_t-3)$) and set
$\alpha=0.1$, $\lambda=0$ and $w_0=0$.

\begin{center}
\begin{tabular}{c|c|c|c}
$t$ & $w_t$ & $g_t=2(w_t-3)$ & $w_{t+1}=w_t-\alpha\,H_t^{-1}g_t$\\ \hline
0 & $0$ & $-6$ & $0-0.1\cdot(1/2)(-6)=0+0.3=0.3$\\
1 & $0.3$ & $2(0.3-3)=-5.4$ & $0.3-0.1\cdot(1/2)(-5.4)=0.3+0.27=0.57$\\
2 & $0.57$ & $2(0.57-3)=-4.86$ & $0.57-0.1\cdot(1/2)(-4.86)=0.57+0.243=0.813$\\
3 & $0.813$ & $2(0.813-3)=-4.374$ & $0.813-0.1\cdot(1/2)(-4.374)=0.813+0.2187=1.0317$
\end{tabular}
\end{center}
Corresponding objective values $f(w_t)$ are $9$, $7.29$, $5.90$, $4.88$, showing monotone decrease.

\newpage
\section*{Assumptions}
\begin{assumption}[Smoothness]\label{A:smooth}
$f$ is $L$‑smooth, i.e.
\[
\forall x,y\in\RR^{d}:\quad 
f(y)\le f(x)+\ip{\nabla f(x)}{y-x}+\frac{L}{2}\norm{y-x}^{2}.
\tag{2}
\]
\end{assumption}
\begin{assumption}[Hessian Lipschitz]\label{A:hesslip}
The Hessian of $f$ is $\rho$‑Lipschitz:
\[
\forall x,y:\;\norm{\nabla^{2}f(x)-\nabla^{2}f(y)}\le\rho\norm{x-y}.
\tag{3}
\]
\end{assumption}
\begin{assumption}[Unbiased stochastic Hessian]\label{A:unbiased}
Conditioned on $x_t$, the minibatch Hessian satisfies
\[
\EE[\,H_t\mid x_t]=\nabla^{2}f(x_t),\qquad 
\EE\bigl[\|H_t-\nabla^{2}f(x_t)\|^{2}\mid x_t\bigr]\le\sigma_{H}^{2}.
\tag{4}
\]
\end{assumption}
\begin{assumption}[Positive definiteness in expectation]\label{A:pd}
There exist constants $0<\mu\le L_H$ such that
\[
\mu I \preceq \nabla^{2}f(x) \preceq L_H I,\qquad\forall x\in\RR^{d}.
\tag{5}
\]
\end{assumption}
\begin{assumption}[Bounded stochastic gradient variance]\label{A:gradvar}
\[
\EE\bigl[\|g_t-\nabla f(x_t)\|^{2}\mid x_t\bigr]\le\sigma_{g}^{2}.
\tag{6}
\]
\end{assumption}

All constants $L,\rho,\mu,L_H,\sigma_H,\sigma_g$ are assumed known; the ridge
parameter $\lambda$ is chosen so that $\tilde H_t\succeq (\mu+\lambda)I$ a.s.

\section*{Main Convergence Theorem}
\begin{theorem}\label{thm:main}
Assume \ref{A:smooth}--\ref{A:gradvar} hold and let the stepsize satisfy
\[
0<\alpha\le\frac{2(\mu+\lambda)}{L\,(\,L_H+\lambda\,)} .
\tag{7}
\]
Then for any $T\ge1$ the iterates generated by \eqref{1} obey
\[
\frac{1}{T}\sum_{t=0}^{T-1}\EE\!\bigl[\|\nabla f(x_t)\|^{2}\bigr]
\;\le\;
\frac{2\bigl(f(x_0)-f^{\star}\bigr)}{\alpha(\mu+\lambda)T}
\;+\;
\frac{L\alpha}{\mu+\lambda}\bigl(\sigma_{g}^{2}+ \tfrac{L_H^{2}}{(\mu+\lambda)^{2}}\sigma_{H}^{2}\bigr),
\tag{8}
\]
where $f^{\star}=\inf_{x}f(x)$.  Consequently,
\[
\min_{0\le t<T}\EE\!\bigl[\|\nabla f(x_t)\|^{2}\bigr]
= O\!\Bigl(\frac{1}{\alpha T}+ \alpha\Bigr).
\]
Choosing $\alpha=\Theta\bigl(T^{-1/2}\bigr)$ yields
\[
\min_{0\le t<T}\EE\!\bigl[\|\nabla f(x_t)\|^{2}\bigr]=O\!\bigl(T^{-1/2}\bigr).
\]
\end{theorem}

\section*{Theoretical Properties}
\begin{itemize}
\item \textbf{Stability.}  The regularisation $\lambda>0$ guarantees that $\tilde H_t$ is uniformly
positive definite, preventing exploding Newton steps.
\item \textbf{Hyper‑parameter sensitivity.}  The admissible stepsize interval \eqref{7}
depends on the ratio $(\mu+\lambda)/(L(L_H+\lambda))$.  Larger minibatches
reduce $\sigma_H,\sigma_g$, allowing a larger $\alpha$ and faster convergence.
\item \textbf{Comparison to SGD.}  If $H_t\approx L I$, the bound \eqref{8} reduces to the classical
$O(1/\sqrt{T})$ rate of SGD, whereas a good curvature estimate shrinks the
prefactor by a factor of $\mu/L$, reflecting the Newton‑type acceleration.
\end{itemize}

\section*{Preliminaries (Definitions and Lemmas)}
\begin{lemma}[Quadratic upper bound from smoothness]\label{lem:smooth}
For any $x$ and any symmetric positive definite $A$,
\[
f\bigl(x-\alpha A^{-1}\nabla f(x)\bigr)
\le f(x)-\alpha\ip{\nabla f(x)}{A^{-1}\nabla f(x)}
+\frac{L\alpha^{2}}{2}\norm{A^{-1}\nabla f(x)}^{2}.
\tag{9}
\]
\end{lemma}
\begin{proof}
Apply the smoothness inequality \eqref{2} with $y=x-\alpha A^{-1}\nabla f(x)$.
\[
\begin{aligned}
f(y) &\le f(x)+\ip{\nabla f(x)}{y-x}
       +\frac{L}{2}\norm{y-x}^{2}\\
     &= f(x)-\alpha\ip{\nabla f(x)}{A^{-1}\nabla f(x)}
        +\frac{L\alpha^{2}}{2}\norm{A^{-1}\nabla f(x)}^{2}.
\end{aligned}
\]
\end{proof}

\begin{lemma}[Matrix norm bound]\label{lem:matnorm}
If $A\succeq (\mu+\lambda)I$, then for any vector $v$,
\[
\frac{1}{L_H+\lambda}\norm{v}^{2}
\;\le\;
v^{\!\top}A^{-1}v
\;\le\;
\frac{1}{\mu+\lambda}\norm{v}^{2}.
\tag{10}
\]
\end{lemma}
\begin{proof}
Since eigenvalues of $A$ lie in $[\mu+\lambda,\,L_H+\lambda]$, the eigenvalues of $A^{-1}$
lie in $[1/(L_H+\lambda),\,1/(\mu+\lambda)]$, which yields \eqref{10}.
\end{proof}

\section*{Main Proof (Step‑by‑Step Derivation)}
We prove Theorem~\ref{thm:main}.  All expectations are taken w.r.t.\ the
randomness of the minibatch $\mathcal S_t$ conditioned on $x_t$ unless otherwise noted.

\medskip
\noindent\textbf{[Step 1] Descent inequality from Lemma~\ref{lem:smooth}.}  
Using $A=\tilde H_t$ in \eqref{9} we obtain
\begin{equation}
f(x_{t+1})
\le f(x_t)-\alpha\ip{\nabla f(x_t)}{\tilde H_t^{-1}\nabla f(x_t)}
     +\frac{L\alpha^{2}}{2}\norm{\tilde H_t^{-1}\nabla f(x_t)}^{2}.
\tag{11}
\end{equation}

\noindent\textbf{[Step 2] Take conditional expectation.}
Condition on $x_t$ and use that $g_t$ is an unbiased estimator of $\nabla f(x_t)$:
\[
\EE[g_t\mid x_t]=\nabla f(x_t),\qquad
\EE[\tilde H_t\mid x_t]=\nabla^{2}f(x_t)+\lambda I .
\]
Because $g_t$ appears only inside $\nabla f(x_t)$ in \eqref{11}, the stochastic gradient
does not affect the first two terms; only the norm term involves $g_t$ when we replace
$\nabla f(x_t)$ by $g_t$ later.  For now we keep $\nabla f(x_t)$.

\noindent\textbf{[Step 3] Bounding the quadratic form.}
From Lemma~\ref{lem:matnorm} and the lower eigenvalue bound $\tilde H_t\succeq(\mu+\lambda)I$,
\[
\ip{\nabla f(x_t)}{\tilde H_t^{-1}\nabla f(x_t)}
\;\ge\;\frac{1}{L_H+\lambda}\norm{\nabla f(x_t)}^{2}.
\tag{12}
\]

\noindent\textbf{[Step 4] Bounding the squared norm term.}
Again by Lemma~\ref{lem:matnorm},
\[
\norm{\tilde H_t^{-1}\nabla f(x_t)}^{2}
\;\le\;\frac{1}{(\mu+\lambda)^{2}}\norm{\nabla f(x_t)}^{2}.
\tag{13}
\]

\noindent\textbf{[Step 5] Insert bounds (12)–(13) into (11).}
\[
\EE\!\bigl[f(x_{t+1})\mid x_t\bigr]
\le f(x_t)-\alpha\frac{1}{L_H+\lambda}\norm{\nabla f(x_t)}^{2}
     +\frac{L\alpha^{2}}{2}\frac{1}{(\mu+\lambda)^{2}}\norm{\nabla f(x_t)}^{2}.
\tag{14}
\]

\noindent\textbf{[Step 6] Isolate the descent coefficient.}
Define
\[
\gamma(\alpha):=
\frac{\alpha}{L_H+\lambda}
-\frac{L\alpha^{2}}{2(\mu+\lambda)^{2}} .
\tag{15}
\]
Then \eqref{14} becomes
\[
\EE\!\bigl[f(x_{t+1})\mid x_t\bigr]
\le f(x_t)-\gamma(\alpha)\,\norm{\nabla f(x_t)}^{2}.
\tag{16}
\]

\noindent\textbf{[Step 7] Choose stepsize to make $\gamma(\alpha)>0$.}
The quadratic $\gamma(\alpha)$ is positive for $0<\alpha<\frac{2(\mu+\lambda)}{L(L_H+\lambda)}$,
which is exactly condition \eqref{7}.  Hence $\gamma(\alpha)\ge
\frac{\alpha(\mu+\lambda)}{2(L_H+\lambda)}$ for any admissible $\alpha$.

\noindent\textbf{[Step 8] Take total expectation and telescope.}
Taking expectation over the history and summing from $t=0$ to $T-1$ yields
\[
\EE[f(x_T)]\le f(x_0)-\gamma(\alpha)\sum_{t=0}^{T-1}\EE\!\bigl[\|\nabla f(x_t)\|^{2}\bigr].
\tag{17}
\]
Since $f(x_T)\ge f^{\star}$,
\[
\sum_{t=0}^{T-1}\EE\!\bigl[\|\nabla f(x_t)\|^{2}\bigr]
\le\frac{f(x_0)-f^{\star}}{\gamma(\alpha)} .
\tag{18}
\]

\noindent\textbf{[Step 9] Replace $\gamma(\alpha)$ by its lower bound.}
Using $\gamma(\alpha)\ge \frac{\alpha(\mu+\lambda)}{2(L_H+\lambda)}$ we obtain
\[
\frac{1}{T}\sum_{t=0}^{T-1}\EE\!\bigl[\|\nabla f(x_t)\|^{2}\bigr]
\le
\frac{2(L_H+\lambda)}{\alpha(\mu+\lambda)}\,
\frac{f(x_0)-f^{\star}}{T}.
\tag{19}
\]

\noindent\textbf{[Step 10] Accounting for stochastic gradient and Hessian noise.}
So far we have ignored the variance terms $\sigma_g^{2}$ and $\sigma_H^{2}$.
To incorporate them we replace the exact gradient $\nabla f(x_t)$ in (11) by the
stochastic gradient $g_t$ and use the decomposition
\[
g_t = \nabla f(x_t) + \varepsilon_t,\qquad
\EE[\varepsilon_t\mid x_t]=0,\;\;
\EE[\|\varepsilon_t\|^{2}\mid x_t]\le\sigma_g^{2}.
\]
Similarly write $H_t = \nabla^{2}f(x_t)+\Delta_t$ with
$\EE[\Delta_t\mid x_t]=0$, $\EE[\|\Delta_t\|^{2}\mid x_t]\le\sigma_H^{2}$.
A straightforward but tedious expansion of
$\norm{\tilde H_t^{-1}g_t}^{2}$ using $(A+B)^{-1}=A^{-1}-A^{-1}BA^{-1}+ \dots$ and
the bounds $\|\Delta_t\|\le\sigma_H$, $\|\varepsilon_t\|\le\sigma_g$ yields
\[
\EE\!\bigl[\norm{\tilde H_t^{-1}g_t}^{2}\mid x_t\bigr]
\le
\frac{1}{(\mu+\lambda)^{2}}\norm{\nabla f(x_t)}^{2}
+ \frac{\sigma_g^{2}}{(\mu+\lambda)^{2}}
+ \frac{L_H^{2}}{(\mu+\lambda)^{4}}\sigma_H^{2}\norm{\nabla f(x_t)}^{2}
+ O(\sigma_H^{2}\sigma_g^{2}).
\tag{20}
\]
Discarding higher‑order mixed terms and inserting (20) into the expectation of
(11) gives an extra additive term
\[
\frac{L\alpha^{2}}{2}\Bigl(
\frac{\sigma_g^{2}}{(\mu+\lambda)^{2}}
+\frac{L_H^{2}}{(\mu+\lambda)^{4}}\sigma_H^{2}\norm{\nabla f(x_t)}^{2}
\Bigr).
\tag{21}
\]
Collecting the coefficients of $\norm{\nabla f(x_t)}^{2}$ we obtain a
modified descent factor
\[
\tilde\gamma(\alpha)=
\frac{\alpha}{L_H+\lambda}
-\frac{L\alpha^{2}}{2(\mu+\lambda)^{2}}
-\frac{L\alpha^{2}L_H^{2}}{2(\mu+\lambda)^{4}}\sigma_H^{2}.
\tag{22}
\]
For the same stepsize range \eqref{7} and a sufficiently small $\sigma_H$,
$\tilde\gamma(\alpha)$ remains positive.  Repeating the telescoping argument
leads to
\[
\frac{1}{T}\sum_{t=0}^{T-1}\EE\!\bigl[\|\nabla f(x_t)\|^{2}\bigr]
\le
\underbrace{\frac{2(L_H+\lambda)}{\alpha(\mu+\lambda)}\frac{f(x_0)-f^{\star}}{T}}_{\text{deterministic term}}
\;+\;
\underbrace{\frac{L\alpha}{\mu+\lambda}\Bigl(\sigma_g^{2}
+ \frac{L_H^{2}}{(\mu+\lambda)^{2}}\sigma_H^{2}\Bigr)}_{\text{variance term}}.
\tag{23}
\]
This is exactly inequality \eqref{8} of Theorem~\ref{thm:main}.

\noindent\textbf{[Step 11] Optimising $\alpha$.}
Balancing the $1/(\alpha T)$ term with the $\alpha$ term in \eqref{23} yields
$\alpha = \Theta(T^{-1/2})$, which gives the $O(T^{-1/2})$ convergence rate
stated after the theorem.

Thus every step required by the theorem has been derived from the
assumptions, completing the proof. \qed

\section*{Rate Derivation (With Constants)}
From \eqref{23} set
\[
C_1:=\frac{2(L_H+\lambda)}{\mu+\lambda}\bigl(f(x_0)-f^{\star}\bigr),\qquad
C_2:=\frac{L}{\mu+\lambda}\Bigl(\sigma_g^{2}+ \frac{L_H^{2}}{(\mu+\lambda)^{2}}\sigma_H^{2}\Bigr).
\]
Then
\[
\frac{1}{T}\sum_{t=0}^{T-1}\EE\!\bigl[\|\nabla f(x_t)\|^{2}\bigr]
\le \frac{C_1}{\alpha T}+C_2\alpha .
\]
Minimising the right‑hand side over $\alpha>0$ gives $\alpha^{\star}
=\sqrt{C_1/(C_2 T)}$, and substituting back yields
\[
\frac{1}{T}\sum_{t=0}^{T-1}\EE\!\bigl[\|\nabla f(x_t)\|^{2}\bigr]
\le 2\sqrt{\frac{C_1 C_2}{T}} = O\!\bigl(T^{-1/2}\bigr).
\]

\section*{Special Cases}
\begin{itemize}
\item \textbf{Exact Hessian, no noise.}  If $\sigma_H=\sigma_g=0$, the bound reduces to
\[
\frac{1}{T}\sum_{t=0}^{T-1}\EE\!\bigl[\|\nabla f(x_t)\|^{2}\bigr]
\le \frac{2(L_H+\lambda)}{\alpha(\mu+\lambda)T}\bigl(f(x_0)-f^{\star}\bigr),
\]
i.e. a pure $O(1/T)$ decay, identical to the deterministic Newton method on a
smooth non‑convex function.
\item \textbf{Large minibatch.}  As $b\to n$, the variances $\sigma_g^{2},\sigma_H^{2}\to0$,
so the algorithm interpolates between SGD ($H_t\approx I$) and the full‑batch
Newton method.
\end{itemize}

\end{document}